% 图 3-1 模块依赖 DAG v4 — 摘要式依赖图：画层级与关键边，不画全量穿越边
\documentclass[border=4pt]{standalone}
\usepackage{xeCJK}
\setCJKmainfont{Songti SC}
\setCJKsansfont{Heiti SC}
\usepackage{tikz}
\usetikzlibrary{arrows.meta, positioning, calc, fit, backgrounds}
% 共享 TikZ 风格 — Aegaeon (SOSP'25) inspired，对标顶刊系统论文
% 使用：% 共享 TikZ 风格 — Aegaeon (SOSP'25) inspired，对标顶刊系统论文
% 使用：% 共享 TikZ 风格 — Aegaeon (SOSP'25) inspired，对标顶刊系统论文
% 使用：\input{_style.tex} 在 standalone preamble 内
% 调色：Okabe-Ito（与 matplotlib 图一致）

\definecolor{fuxiBlue}{HTML}{0072B2}    % 自家系统主色
\definecolor{fuxiOrange}{HTML}{E69F00}  % 第一对照
\definecolor{fuxiGreen}{HTML}{009E73}
\definecolor{fuxiPurple}{HTML}{CC79A7}
\definecolor{fuxiSky}{HTML}{56B4E9}
\definecolor{fuxiGrayBox}{HTML}{F2F2F2} % 浅灰底
\definecolor{fuxiGrayLine}{HTML}{808080}
\definecolor{fuxiGrayDark}{HTML}{555555}

\tikzset{
  % ── 通用节点 ────────────────────────────────────────
  modBox/.style = {
    draw=fuxiGrayLine, thick, fill=white,
    rounded corners=2pt, inner sep=4pt,
    minimum height=8mm, minimum width=20mm,
    align=center, font=\small
  },
  highlightBox/.style = {
    modBox, draw=fuxiBlue, fill=fuxiBlue!8,
    text=fuxiBlue, font=\small\bfseries
  },
  ghostBox/.style = {
    modBox, draw=fuxiGrayLine!50, fill=fuxiGrayBox,
    text=fuxiGrayDark, font=\small\itshape
  },
  layerBg/.style = {
    draw=fuxiGrayLine!40, dashed, fill=fuxiGrayBox!40,
    rounded corners=4pt, inner sep=8pt
  },
  externalBoundary/.style = {
    draw=fuxiGrayLine, dashed, dash pattern=on 4pt off 3pt,
    rounded corners=6pt, inner sep=10pt
  },
  % ── 箭头 ─────────────────────────────────────────────
  flow/.style = {
    -{Stealth[length=4pt, width=4pt]},
    line width=0.7pt, color=fuxiGrayDark
  },
  flowEvent/.style = {
    flow, color=fuxiBlue, line width=0.9pt
  },
  flowDup/.style = {
    flow, color=fuxiBlue!70, line width=0.7pt, dotted
  },
  bidir/.style = {
    {Stealth[length=4pt]}-{Stealth[length=4pt]},
    line width=0.7pt, color=fuxiGrayDark
  },
  % ── 编号气泡 ① ② ③ ────────────────────────────────
  numBubble/.style = {
    circle, draw=fuxiBlue, fill=white, text=fuxiBlue,
    font=\scriptsize\bfseries, inner sep=0pt,
    minimum size=4mm, line width=0.6pt
  },
}
 在 standalone preamble 内
% 调色：Okabe-Ito（与 matplotlib 图一致）

\definecolor{fuxiBlue}{HTML}{0072B2}    % 自家系统主色
\definecolor{fuxiOrange}{HTML}{E69F00}  % 第一对照
\definecolor{fuxiGreen}{HTML}{009E73}
\definecolor{fuxiPurple}{HTML}{CC79A7}
\definecolor{fuxiSky}{HTML}{56B4E9}
\definecolor{fuxiGrayBox}{HTML}{F2F2F2} % 浅灰底
\definecolor{fuxiGrayLine}{HTML}{808080}
\definecolor{fuxiGrayDark}{HTML}{555555}

\tikzset{
  % ── 通用节点 ────────────────────────────────────────
  modBox/.style = {
    draw=fuxiGrayLine, thick, fill=white,
    rounded corners=2pt, inner sep=4pt,
    minimum height=8mm, minimum width=20mm,
    align=center, font=\small
  },
  highlightBox/.style = {
    modBox, draw=fuxiBlue, fill=fuxiBlue!8,
    text=fuxiBlue, font=\small\bfseries
  },
  ghostBox/.style = {
    modBox, draw=fuxiGrayLine!50, fill=fuxiGrayBox,
    text=fuxiGrayDark, font=\small\itshape
  },
  layerBg/.style = {
    draw=fuxiGrayLine!40, dashed, fill=fuxiGrayBox!40,
    rounded corners=4pt, inner sep=8pt
  },
  externalBoundary/.style = {
    draw=fuxiGrayLine, dashed, dash pattern=on 4pt off 3pt,
    rounded corners=6pt, inner sep=10pt
  },
  % ── 箭头 ─────────────────────────────────────────────
  flow/.style = {
    -{Stealth[length=4pt, width=4pt]},
    line width=0.7pt, color=fuxiGrayDark
  },
  flowEvent/.style = {
    flow, color=fuxiBlue, line width=0.9pt
  },
  flowDup/.style = {
    flow, color=fuxiBlue!70, line width=0.7pt, dotted
  },
  bidir/.style = {
    {Stealth[length=4pt]}-{Stealth[length=4pt]},
    line width=0.7pt, color=fuxiGrayDark
  },
  % ── 编号气泡 ① ② ③ ────────────────────────────────
  numBubble/.style = {
    circle, draw=fuxiBlue, fill=white, text=fuxiBlue,
    font=\scriptsize\bfseries, inner sep=0pt,
    minimum size=4mm, line width=0.6pt
  },
}
 在 standalone preamble 内
% 调色：Okabe-Ito（与 matplotlib 图一致）

\definecolor{fuxiBlue}{HTML}{0072B2}    % 自家系统主色
\definecolor{fuxiOrange}{HTML}{E69F00}  % 第一对照
\definecolor{fuxiGreen}{HTML}{009E73}
\definecolor{fuxiPurple}{HTML}{CC79A7}
\definecolor{fuxiSky}{HTML}{56B4E9}
\definecolor{fuxiGrayBox}{HTML}{F2F2F2} % 浅灰底
\definecolor{fuxiGrayLine}{HTML}{808080}
\definecolor{fuxiGrayDark}{HTML}{555555}

\tikzset{
  % ── 通用节点 ────────────────────────────────────────
  modBox/.style = {
    draw=fuxiGrayLine, thick, fill=white,
    rounded corners=2pt, inner sep=4pt,
    minimum height=8mm, minimum width=20mm,
    align=center, font=\small
  },
  highlightBox/.style = {
    modBox, draw=fuxiBlue, fill=fuxiBlue!8,
    text=fuxiBlue, font=\small\bfseries
  },
  ghostBox/.style = {
    modBox, draw=fuxiGrayLine!50, fill=fuxiGrayBox,
    text=fuxiGrayDark, font=\small\itshape
  },
  layerBg/.style = {
    draw=fuxiGrayLine!40, dashed, fill=fuxiGrayBox!40,
    rounded corners=4pt, inner sep=8pt
  },
  externalBoundary/.style = {
    draw=fuxiGrayLine, dashed, dash pattern=on 4pt off 3pt,
    rounded corners=6pt, inner sep=10pt
  },
  % ── 箭头 ─────────────────────────────────────────────
  flow/.style = {
    -{Stealth[length=4pt, width=4pt]},
    line width=0.7pt, color=fuxiGrayDark
  },
  flowEvent/.style = {
    flow, color=fuxiBlue, line width=0.9pt
  },
  flowDup/.style = {
    flow, color=fuxiBlue!70, line width=0.7pt, dotted
  },
  bidir/.style = {
    {Stealth[length=4pt]}-{Stealth[length=4pt]},
    line width=0.7pt, color=fuxiGrayDark
  },
  % ── 编号气泡 ① ② ③ ────────────────────────────────
  numBubble/.style = {
    circle, draw=fuxiBlue, fill=white, text=fuxiBlue,
    font=\scriptsize\bfseries, inner sep=0pt,
    minimum size=4mm, line width=0.6pt
  },
}


\begin{document}
\begin{tikzpicture}[font=\sffamily, node distance=8mm and 10mm]

% ── Layer rows ─────────────────────────────────────
% 先放 L4 行（cc 在原点），cli 再 above=cc 自动居中
\node[modBox, minimum width=31mm, minimum height=10mm] (cc)
  {fuxi-agent-cc};
\node[modBox, left=8mm of cc, minimum width=31mm, minimum height=10mm] (codex)
  {fuxi-agent-codex};
\node[modBox, right=8mm of cc, minimum width=31mm, minimum height=10mm] (sched)
  {fuxi-scheduler\\\scriptsize cron};

% 显式 at(cc.north + 上方 10mm) 并锚 south, 强制 cli 中心 x = cc 中心 x
\node[highlightBox, anchor=south, minimum width=36mm, minimum height=11mm] (cli)
  at ($(cc.north)+(0,10mm)$) {fuxi-cli\\\scriptsize daemon + binaries};

\node[highlightBox, below=12mm of cc, minimum width=36mm, minimum height=11mm] (orch)
  {fuxi-orchestrator\\\scriptsize Fuxi + Shelf};
\node[modBox, left=10mm of orch, minimum width=31mm, minimum height=10mm] (firehose)
  {fuxi-firehose\\\scriptsize WS hub + TUI};
\node[modBox, right=10mm of orch, minimum width=31mm, minimum height=10mm] (im)
  {fuxi-im\\\scriptsize PWA 后端};

\node[modBox, below left=12mm and 12mm of orch, minimum width=32mm, minimum height=10mm] (events)
  {fuxi-events\\\scriptsize EventBus + WAL};
\node[modBox, below=12mm of orch, minimum width=32mm, minimum height=10mm] (memory)
  {fuxi-memory\\\scriptsize 抽取/召回};
\node[modBox, below right=12mm and 12mm of orch, minimum width=32mm, minimum height=10mm] (a2a)
  {fuxi-a2a\\\scriptsize wire + SSE};

\node[highlightBox, below=14mm of memory, minimum width=38mm, minimum height=11mm] (core)
  {fuxi-core\\\scriptsize trait + 基础类型};

% ── Layer backgrounds ──────────────────────────────
\begin{scope}[on background layer]
  \node[layerBg, fit=(cli), inner sep=5pt,
    label={[anchor=east, font=\scriptsize\itshape, text=fuxiGrayDark]left:{L5 顶层}}] {};
  \node[layerBg, fit=(codex) (cc) (sched), inner sep=5pt,
    label={[anchor=east, font=\scriptsize\itshape, text=fuxiGrayDark]left:{L4 适配/触发}}] {};
  \node[layerBg, fit=(firehose) (orch) (im), inner sep=5pt,
    label={[anchor=east, font=\scriptsize\itshape, text=fuxiGrayDark]left:{L3 编排/观测}}] {};
  \node[layerBg, fit=(events) (memory) (a2a), inner sep=5pt,
    label={[anchor=east, font=\scriptsize\itshape, text=fuxiGrayDark]left:{L2 基础设施}}] {};
  \node[layerBg, fit=(core), inner sep=5pt,
    label={[anchor=east, font=\scriptsize\itshape, text=fuxiGrayDark]left:{L1 核心}}] {};
\end{scope}

% ── Dependency edges：树形布局 — trunk + rail + 单箭头进中心,不再贴边 ──

% L5 → L4：fan-out 树（cli 单干 → 横 rail → 3 条下垂分支）
\coordinate (cliRailY) at ($(cli.south)+(0,-4mm)$);
\draw[color=fuxiGrayDark, line width=0.7pt] (cli.south) -- (cliRailY);
\draw[color=fuxiGrayDark, line width=0.7pt] (codex.north |- cliRailY) -- (sched.north |- cliRailY);
\draw[flow] (codex.north |- cliRailY) -- (codex.north);
\draw[flow] (cc.north    |- cliRailY) -- (cc.north);
\draw[flow] (sched.north |- cliRailY) -- (sched.north);

% L4 → L3：fan-in 树（3 条上升 stem → 横 rail → 单箭头进 orch.north 中心）
\coordinate (orchRailY) at ($(orch.north)+(0,4mm)$);
\draw[color=fuxiGrayDark, line width=0.7pt] (codex.south) -- (codex.south |- orchRailY);
\draw[color=fuxiGrayDark, line width=0.7pt] (cc.south)    -- (cc.south    |- orchRailY);
\draw[color=fuxiGrayDark, line width=0.7pt] (sched.south) -- (sched.south |- orchRailY);
\draw[color=fuxiGrayDark, line width=0.7pt] (codex.south |- orchRailY) -- (sched.south |- orchRailY);
\draw[flow] (orch.north |- orchRailY) -- (orch.north);

% L3 同层（peer）
\draw[flow] (orch.west) -- (firehose.east);
\draw[flow] (orch.east) -- (im.west);

% L3 → L2：fan-out 树（orch 单干 → 横 rail → 3 条下垂分支）
\coordinate (orchOutY) at ($(orch.south)+(0,-4mm)$);
\draw[color=fuxiGrayDark, line width=0.7pt] (orch.south) -- (orchOutY);
\draw[color=fuxiGrayDark, line width=0.7pt] (events.north |- orchOutY) -- (a2a.north |- orchOutY);
\draw[flow] (events.north |- orchOutY) -- (events.north);
\draw[flow] (memory.north |- orchOutY) -- (memory.north);
\draw[flow] (a2a.north    |- orchOutY) -- (a2a.north);

% L2 → L1：fan-in 树（3 条上升 stem → 横 rail → 单箭头进 core.north 中心）
\coordinate (coreRailY) at ($(core.north)+(0,4mm)$);
\draw[color=fuxiGrayDark, line width=0.7pt] (events.south) -- (events.south |- coreRailY);
\draw[color=fuxiGrayDark, line width=0.7pt] (memory.south) -- (memory.south |- coreRailY);
\draw[color=fuxiGrayDark, line width=0.7pt] (a2a.south)    -- (a2a.south    |- coreRailY);
\draw[color=fuxiGrayDark, line width=0.7pt] (events.south |- coreRailY) -- (a2a.south |- coreRailY);
\draw[flow] (core.north |- coreRailY) -- (core.north);

% ── Key invariant note ──────────────────────────────
\node[anchor=north, font=\scriptsize\itshape, text=fuxiOrange, align=center,
      draw=fuxiOrange!40, fill=fuxiOrange!6, rounded corners=2pt, inner sep=4pt]
  at ($(core.south)+(0,-8mm)$) {%
    关键约束：\texttt{fuxi-memory}\,$\not\to$\,\texttt{fuxi-orchestrator}；
    adapter 实现统一放在 \texttt{fuxi-cli} 顶层
  };

\end{tikzpicture}
\end{document}
